%%%%%%%%%%  scan idx 587  |  folio 574  %%%%%%%%%%
\SGAPrintPage{574}%
consequence, the immersions \(x_i:S\to C\) (\(i=0,1\)) are regular and remain regular immersions \(x_i:s\to C_s\) under the base change \(s\to S\). Note that, by the exact homotopy sequence (IX 1.1), \(E_t\in K_\bullet(X_t\times_{k(t)}C_t)\) comes from an element \(E\in K_\bullet(X\times_SC)\), so the element \(E_t(x_{1t})-E_t(x_{0t})\) of \(K_\bullet(X_t)\) comes from the element \(E(x_1)-E(x_0)\) of \(K_\bullet(X\times_SC)\), and consequently its image in \(K_\bullet(X_s)\) under (7.3.1) is precisely
\sgaStableTarget{sga6.e10.eq.unnamed.occ2}\[
  i^*(E(x_1)-E(x_0)).
  \tag{*}
\]
Of course, for a section \(x\) of \(C\) over \(S\), one puts
\[
  E(x)=i_x^*(E),\qquad\text{where } i_x:X\longrightarrow X\times_SC
\]
is defined as above, and can be regarded as obtained from \(x\) by the flat base change \(X\times_SC\to C\); it is therefore a regular immersion, a fortiori a perfect morphism, which gives a meaning to the expression \(E(x)\). For the same reason, the expression \(E_s(x_s)\) is defined, and, by transitivity of inverse images (IV 2.12) in the following commutative diagram of perfect morphisms:
\[
\begin{tikzcd}[column sep=large,row sep=large]
X \arrow[r,"i_x"] & X\times_SC\\
X_s \arrow[u,"i"] \arrow[r,"i_{x_s}"] & X_s\times_{k(s)}C_s \arrow[u]
\end{tikzcd}
\;,
\]
one finds that the left-hand side of \((*)\) is equal to
\sgaStableTarget{sga6.e10.eq.unnamed.occ3}\[
  E_s(x_{1s})-E_s(x_{0s}).
  \tag{**}
\]
Since \(k(s)\) is algebraically closed, the reduced irreducible components of \(C_s\) are geometrically integral over \(k(s)\) and their intersection points are rational over \(k(s)\). One concludes readily that the expression \((**)\) is algebraically equivalent to zero in \(K_\bullet(X_s)\), using the connectedness

%%%%%%%%%%  scan idx 588  |  folio 575  %%%%%%%%%%
\SGAPrintPage{575}%

of \(C_s\), which follows from the geometric connectedness of \(C_t\).

From the homomorphism (7.12.1), using 7.4, one obtains a homomorphism
\sgaStableTarget{sga6.e10.eq.7-12-2}\begin{equation}\tag{7.12.2}
  \Gr_\bullet^{\mathrm{alg}}(X_t)\longrightarrow \Gr_\bullet^{\mathrm{alg}}(X_s).
\end{equation}

When \(X_s\) and \(X_t\) are moreover assumed regular, (7.12.1) can also be interpreted as a ring homomorphism
\sgaStableTarget{sga6.e10.eq.7-12-3}\begin{equation}\tag{7.12.3}
  K^\bullet_{\mathrm{alg}}(X_t)\longrightarrow K^\bullet_{\mathrm{alg}}(X_s),
\end{equation}
which, when \(X\) is moreover separated, is a homomorphism of \(\lambda\)-rings, and consequently defines a homomorphism
\sgaStableTarget{sga6.e10.eq.7-12-4}\begin{equation}\tag{7.12.4}
  \Gr^\bullet_{\mathrm{alg}}(X_t)\longrightarrow \Gr^\bullet_{\mathrm{alg}}(X_s).
\end{equation}

\sgaStableTarget{sga6.e10.para.7-12-5}\textbf{7.12.5.} Let \(T\) be a smooth proper algebraic scheme over a field \(k\). One checks easily that an element \(x\) of \(K^\bullet(T)\), for which there is an integer \(n\geq1\) with \(nx\) algebraically equivalent to zero, is numerically equivalent to zero. It is plausible that the converse is always true, at least when \(T\) is projective. In this question, one can plainly suppose that \(k\) is algebraically closed, and replace \(K^\bullet(T)\) by \(\Gr^\bullet(T)\) or \(\Gr_\bullet(T)\) without changing the conjecture. One also obtains an equivalent conjecture by working with the Chow ring \(A(T)\) (cf. XIV 4.1), obtaining the following more or less classical conjecture: if \(Z\) is an algebraic cycle numerically equivalent to zero on \(T\), there is an integer \(n\geq1\) with \(nZ\) algebraically equivalent to zero. In the case where this conjecture is correct for \(X_t\), the compatibility of the homomorphism (7.3.1) with algebraic equivalence would imply that this specialization homomorphism (when \(f:X\to S\) is a projective smooth morphism) is also compatible with numerical equivalence, as conjectured in 7.9.

%%%%%%%%%%  scan idx 589  |  folio 576  %%%%%%%%%%
\SGAPrintPage{576}%

\sgaStableTarget{sga6.e10.sec.7.13}\subsection*{7.13. Relations with \(\ell\)-adic cohomology}
Recall (SGA 5 VII) that, for every locally free Module of finite type \(\underline{E}\) on a scheme \(T\), and every prime number distinct from the residue characteristics of \(T\), one defines \(\ell\)-adic Chern classes (for \(i\geq 0\))
\[
  c^i(E)\in H^{2i}(T,\mathbb Z_\ell(i)),
\]
satisfying an additivity formula
\[
  c(E)=c(E')c(E'')
\]
when \(E\) is an extension of a locally free module \(E'\) by a locally free module \(E''\). From this additivity one concludes that the total Chern class \(c(E)=\sum_i c^i(E)\) arises from a homomorphism
\sgaStableTarget{sga6.e10.eq.7-13-0}\begin{equation}\tag{7.13.0}
  c:K^\bullet(T)\longrightarrow
  \widehat G\bigl(H^{2*}(T,\mathbb Z_\ell(*))\bigr)
  =
  1+\prod_{i\geq 1}H^{2i}(T,\mathbb Z_\ell(i)),
\end{equation}
(where the target is a group under multiplication, as explained in V 1, and where \(K^\bullet(T)\) is the ``\emph{naive}'' \(K^\bullet\)). One obtains canonical maps
\sgaStableTarget{sga6.e10.eq.7-13-1}\begin{equation}\tag{7.13.1}
  c^i:K^\bullet(T)\longrightarrow H^{2i}(T,\mathbb Z_\ell(i)).
\end{equation}

Now let \(X\) be a \emph{proper} flat scheme over \(S\), with \emph{regular fibers}, and suppose that \(S\) is \emph{henselian}. Assume that \(\ell\neq\) the residue characteristic of \(S\). One can then form the diagram
\sgaStableTarget{sga6.e10.eq.7-13-2}\begin{equation}\tag{7.13.2}
\begin{tikzcd}[column sep=large,row sep=large]
K^\bullet(X_t) \arrow[r,"c^i"] \arrow[d,"\sigma"'] &
H^{2i}(X_t,\mathbb Z_\ell(i))\\
K^\bullet(X_s) \arrow[r,"c^i"] & H^{2i}(X_s,\mathbb Z_\ell(i)) \arrow[u,"\sigma^{-1}"']
\end{tikzcd}
\;,
\end{equation}
where the horizontal arrows are the ``Chern class'' maps for \(X_t\) and \(X_s\) respectively, and where \(\sigma^{-1}\) is the composite homomorphism
\[
  H(X_s)\xleftarrow{\substack{i^*\\\sim}}H(X)\xrightarrow{j^*}H(X_t);
\]

%%%%%%%%%%  scan idx 590  |  folio 577  %%%%%%%%%%
\SGAPrintPage{577}%

the arrows displayed here are restriction homomorphisms, and one uses the fact that \(i^*\) is an isomorphism by the proper base change theorem (SGA 4 XII 5.5). I claim that \emph{the square} (7.13.2) \emph{is commutative}. The proof is essentially the same as that establishing the commutativity of (7.10.2); here it suffices to use the compatibility of the formation of Chern classes with inverse images for the morphisms \(j:X_t\to X\) and \(i:X_s\to X\).

\sgaStableTarget{sga6.e10.para.7-13-3}\textbf{7.13.3.} Note that the square (7.13.2) arises by passage to the projective limit from the analogous squares in which cohomology with coefficients in the \(\ell\)-adic sheaf \(\mathbb Z_\ell(i)\) is replaced by cohomology with coefficients in the torsion sheaves \(\mu_n^{\otimes i}\) (where \(n\) runs through the powers of \(\ell\)). Thus the commutativity of (7.13.2) merely expresses the commutativity of the squares involving Chern classes mod. \(n\), for every power \(n\) of \(\ell\). Choose an algebraic closure \(\bar K\) of \(K=k(t)\), and hence an algebraic closure \(\bar k\) of \(k=k(s)\) (namely the residue field of the integral closure of \(V\) in \(\bar K\)). For every finite subextension \(K_\alpha\) of \(\bar K\), one can consider the analogous commutative diagram (for fixed \(n\)) in the situation obtained from \(f:X\to S\) by the base change \(S_\alpha\to S\), with \(S_\alpha\) the normalization of \(S\) in \(K_\alpha\). For this, note that, by Krull--Akizuki and since \(S\) is henselian, \(S_\alpha\) is indeed a trait; moreover, one must suppose here that the fibers of \(X\) over \(S\) are geometrically regular, i.e., that \(X\) is \emph{smooth} over \(S\), so that this hypothesis is preserved by base change. As \(\alpha\) varies, one thus obtains an inductive system of commutative squares. Passing to the inductive limit and using (SGA 4 VII 5.8) and (IV 3.2), one obtains a commutative square (cf. 7.16):

%%%%%%%%%%  scan idx 591  |  folio 578  %%%%%%%%%%
\SGAPrintPage{578}%

\[
\begin{tikzcd}[column sep=large,row sep=large]
K^\bullet(X_{\bar t}) \arrow[r,"c^i"] \arrow[d,"\sigma"'] &
H^{2i}(X_{\bar t},\mu_n^{\otimes i})\\
K^\bullet(X_{\bar s}) \arrow[r,"c^i"] &
H^{2i}(X_{\bar s},\mu_n^{\otimes i}) \arrow[u,"\sigma^{-1}"]
\end{tikzcd}
\;,
\]
where \(X_{\bar t}\), \(X_{\bar s}\) are the \emph{geometric} fibers of \(X\). Moreover, one knows that the cohomological specialization homomorphism \(\sigma^{-1}\) is now an isomorphism (SGA 4 XVI 2.2), so the commutativity of the preceding square can also be expressed as the commutativity of the square obtained by reversing the preceding arrow. Passing to the projective limit over \(n\) in the commutative squares thus obtained, one finds a commutative square
\sgaStableTarget{sga6.e10.eq.7-13-4}\begin{equation}\tag{7.13.4}
\begin{tikzcd}[column sep=large,row sep=large]
K^\bullet(X_{\bar t}) \arrow[r] \arrow[d,"\sigma"'] &
H^{2i}(X_{\bar t},\mathbb Z_\ell(i)) \arrow[d,"\wr"']\\
K^\bullet(X_{\bar s}) \arrow[r] &
H^{2i}(X_{\bar s},\mathbb Z_\ell(i)).
\end{tikzcd}
\end{equation}

\sgaStableTarget{sga6.e10.para.7-13-5}\textbf{7.13.5.} Another way to express this commutativity is to say that, when \(x_t\) is an element of \(K^\bullet(X_t)\) and \(x_s\) is its image in \(K^\bullet(X_s)\), \(c^i(x_{\bar t})\) and \(c^i(x_{\bar s})\) are the values, at the geometric points \(\bar t\) and \(\bar s\) respectively, of a \emph{section} of \(R^{2i}f_*(\mathbb Z_\ell(i))\). This section is moreover readily interpreted by taking an element \(x\in K^\bullet(X)\) inducing \(x_t\), whence a \(c^i(x)\in H^{2i}(X,\mathbb Z_\ell(i))\), which defines the required section of \(R^{2i}f_*(\mathbb Z_\ell(i))\).

\sgaStableTarget{sga6.e10.para.7-13-6}\textbf{7.13.6.} Let \(T\) be a scheme. Then, by passage to the associated graded groups, the homomorphism (7.13.0) defines homomorphisms
\[
  \psi_i:\Gr^i(T)\longrightarrow H^{2i}(T,\mathbb Z_\ell(i)),
\]
the image under \(\psi_i\) of an element \(x\) of \(\Filt^iK^\bullet(T)\) being, by definition, \(c^i(x)\). (The fact that (7.13.0) is compatible with the \(\lambda\)-filtration on
